\subsectionaddtolist{الف}

سیگنال داده‌شده را می‌توان به صورت مجموع دو سیگنال نوشت:
\[
x[n] = n \left(\frac{1}{2}\right)^{|n|} = n \left(\frac{1}{2}\right)^n u[n] + n \left(\frac{1}{2}\right)^{-n} u[-n-1]
\]
\[
x[n] = n \left(\frac{1}{2}\right)^n u[n] + n (2)^n u[-n-1]
\]

با استفاده از خاصیت مشتق‌گیری در حوزه Z که به صورت $\mathcal{Z}\{n g[n]\} = -z \frac{dG(z)}{dz}$ تعریف می‌شود، تبدیل Z هر بخش را به دست می‌آوریم:
\begin{enumerate}
    \item برای بخش اول $x_1[n] = n \left(\frac{1}{2}\right)^n u[n]$:
    \[
    \mathcal{Z}\left\{\left(\frac{1}{2}\right)^n u[n]\right\} = \frac{1}{1 - \frac{1}{2}z^{-1}} \quad \Rightarrow \quad X_1(z) = -z \frac{d}{dz} \left( \frac{1}{1 - \frac{1}{2}z^{-1}} \right) = \frac{\frac{1}{2}z^{-1}}{\left(1 - \frac{1}{2}z^{-1}\right)^2}
    \]
    
    \item برای بخش دوم $x_2[n] = n (2)^n u[-n-1]$:
    \[
    \mathcal{Z}\left\{(2)^n u[-n-1]\right\} = -\frac{1}{1 - 2z^{-1}} \quad \Rightarrow \quad X_2(z) = -z \frac{d}{dz} \left( -\frac{1}{1 - 2z^{-1}} \right) = \frac{-2z^{-1}}{\left(1 - 2z^{-1}\right)^2}
    \]
\end{enumerate}

بنابراین، تبدیل Z کل سیگنال در فرم بسته برابر است با:
\[
X(z) = \frac{\frac{1}{2}z^{-1}}{\left(1 - \frac{1}{2}z^{-1}\right)^2} - \frac{2z^{-1}}{\left(1 - 2z^{-1}\right)^2}
\]

\subsectionaddtolist{ب}

ناحیه همگرایی کل، اشتراک نواحی همگرایی دو بخش فوق است:
\begin{itemize}
    \item برای جمله اول: $|z| > \frac{1}{2}$
    \item برای جمله دوم: $|z| < 2$
\end{itemize}
با اشتراک گرفتن از این دو بازه، ناحیه همگرایی به دست می‌آید:
\[
\frac{1}{2} < |z| < 2
\]

\subsectionaddtolist{ج}

یک سیگنال زمان-گسسته دارای تبدیل فوریه (DTFT) است اگر و تنها اگر ناحیه همگرایی تبدیل Z آن شامل دایره واحد ($|z| = 1$) باشد.
\newline
با توجه به ناحیه همگرایی به دست آمده در بخش قبل ($\frac{1}{2} < |z| < 2$)، دایره واحد ($|z| = 1$) کاملا درون این ناحیه قرار دارد. بنابراین، این سیگنال {دارای تبدیل فوریه زمان-گسسته است}.
